\chapter{Historical and Archival Methods}
\label{ch:historical-archival}

\epigraph{History is the memory of states.}{Henry Kissinger}

\noindent
Doxonomic systems operate over decades and centuries.
Computational methods capture the present; historical methods reveal how we arrived here.
This chapter develops the tools for doxonomic history: genealogy of concepts, institutional archives, long-duration narrative change, and the integration of historical and computational evidence.

\section{Genealogy of Concepts}
\label{sec:ch18-genealogy}

\begin{definition}[Doxonomic Genealogy]
\label{def:genealogy}
\index{genealogy}
A \emph{doxonomic genealogy} of belief $p$ is a reconstruction of:
\begin{enumerate}[nosep]
  \item when $p$ was first articulated as an explicit claim,
  \item which agents promoted it and with what resources,
  \item through which institutional channels it gained credibility,
  \item what alternative beliefs it displaced,
  \item at what point it achieved common-sense status (if ever).
\end{enumerate}
\end{definition}

This extends \textcite{foucault1971} from discourse analysis to funded discourse analysis: not just \emph{who said what when}, but \emph{who paid for it to be said, through what institutions, with what effect}.

\section{Institutional Archives}
\label{sec:ch18-archives}

Key archival sources for doxonomic history:
\begin{itemize}[nosep]
  \item foundation annual reports and grant records,
  \item corporate archives (advertising, internal strategy),
  \item government propaganda agency records,
  \item educational curriculum archives,
  \item media ownership transfer records,
  \item think-tank publication archives.
\end{itemize}

\section{Long-Duration Narrative Change}
\label{sec:ch18-longue-duree}

\begin{model}[Periodization of Belief Regimes]
\label{mod:periodization}
\index{periodization}
A doxonomic period is defined by the dominance of an anthropological regime $\alpha$.
Transitions between periods are marked by:
\begin{itemize}[nosep]
  \item shifts in funding patterns (new donors, new institutional vehicles),
  \item institutional realignment (new think tanks, curriculum reforms),
  \item narrative crises (events that destabilize the dominant frame),
  \item measurable changes in population belief distributions.
\end{itemize}
\end{model}

\begin{example}
The shift from Keynesian common sense (``government manages the economy'') to neoliberal common sense (``markets are self-correcting'') can be periodized as: intellectual groundwork (1947--1970, Mont P\`{e}lerin Society), institutional build-out (1970--1980, Heritage Foundation, Cato Institute, AEI), political capture (1979--1983, Thatcher and Reagan), and normalization (1990--2008, Third Way, Washington Consensus) \autocite{mirowski2009, harvey2005}.
\end{example}

\section{Notes and References}
\label{sec:ch18-notes}

Genealogical method draws on \textcite{foucault1971}.
The institutional history of neoliberalism is documented in \textcite{mirowski2009} and \textcite{mayer2016}.
On long-duration institutional change, see \textcite{north1990} and \textcite{acemoglu2012}.
Archival methods for social science are covered in \textcite{king1994}.

\section{Exercises}

\begin{exercise}
Construct a doxonomic genealogy of the belief ``taxation is theft'' from its philosophical origins to its current institutional support.
\end{exercise}

\begin{exercise}
Using the periodization framework, identify the major doxonomic periods for beliefs about gender roles in the 20th century.
What institutional shifts mark the transitions?
\end{exercise}

\begin{exercise}
Design an archival research project to estimate the total spending on neoliberal narrative production from 1970 to 2000.
What sources would you use?
What are the main obstacles?
\end{exercise}
